% !TEX program = pdflatex
% !TEX encoding = UTF-8 Unicode

% Plantilla, basada en la clase `scrbook` del paquete KOMA-script,  para la elaboración de un TFG siguiendo las directrices del la comisión del Grado en Matemáticas de la Universidad de Granada.

% Francisco Torralbo Torralbo

\documentclass[print, color]{ugrTFG}

% VERSIÓN ELECTRÓNICA PARA TABLETA
% Cambiando la opción "print" por "tablet" generaremos un pdf adaptado para leerlo en tabletas de 9 pulgadas.

\usepackage{amstext} % for \text macro
\usepackage{array}   % for \newcolumntype macro
\usepackage[spanish]{babel}
\newcolumntype{L}{>{$}l<{$}} % math-mode version of "l" column type
\newcolumntype{C}{>{$}c<{$}} % math-mode version of "l" column type

\newcommand{\R}{\mathbb{R}}
\newcommand{\MAS}{\text{Muestra Aleatoria Simple}}
% \newcommand{\P}{\mathbb{P}}

% -------------------------------------------------------------------
% INFORMACIÓN DEL TFG Y EL AUTOR
% -------------------------------------------------------------------

\newcommand{\miTitulo}{Multivariable regression versus the G-Formula\xspace}
\newcommand{\miNombre}{Lucas Fehlau Arbulu\xspace}
\newcommand{\miGrado}{Grado en Matemáticas}
\newcommand{\miFacultad}{Facultad de Ciencias}
\newcommand{\miUniversidad}{Universidad de Granada}

% Añadir tantos tutores como sea necesario separando cada uno de ellos mediante el comando `\medskip` y una línea en blanco
\newcommand{\miTutor}{
  Miguel Ángel Luque Fernández \\ \emph{Departamento de Estadística e Investigación Operativa}

  % Añadir tantos tutores como sea necesario.

  % \medskip
   %Nombre del tutor 2 \\ \emph{Departamento del tutor 2}
}
\newcommand{\miCurso}{2025-2026\xspace}

\hypersetup{
  pdftitle={\miTitulo},
  pdfauthor={\textcopyright\ \miNombre, \miFacultad, \miUniversidad}
}

\begin{document}

¿
\textquestiondown
% -------------------------------------------------------------------
% GLOSARIO: Opcional
% -------------------------------------------------------------------

% !TeX root = ../tfg.tex
% !TeX encoding = utf8

\chapter*{Glosario}
\addcontentsline{toc}{chapter}{Glosario} % Añade el glosario a la tabla de contenidos

\section*{Notación}
\begin{description}
    \item[$\R$] Conjunto de números reales.
    \item[$\N$] Conjunto de números naturales.
    \item[$\E$] Esperanza de una variable aleatoria.
    \item[$\P$] Probabilidad.
    \item[$A$] Tratamiento o exposición, binario salvo mención de lo contrario.
    \item[$Y$] Resultado o respuesta (\textit{outcome}).
    \item[$X$] Covariantes.
    \item[$Y(a)$] Resultado potencial bajo el tratamiento $A = a$.
    \item[$\cdot \indep \cdot$] Independencia de variables aleatorias.
    \item[$\cdot \indep \cdot \, | \, \cdot $] Independencia condicional de variables aleatorias.
    \item[$\mu_a(x)$] Esperanza condicional $\E[Y | A = a, X = x]$.
    \item[$e(x)$] Propensidad al tratamiento, $P(A = 1 | X = x)$.
    \item[$PA_i$] Padres markovianos de $X_i$.
    \item[$\hat \theta$] Estimador de la cantidad $\theta$.
\end{description}

\section*{Siglas}
\begin{description}
    \item[IID] Independientes e idénticamente distribuidas.
    \item[MAS] Muestra aleatoria simple.
    \item[RCT] \textit{Randomized Controlled Trial}, estudio controlado aleatorizado.
    \item[SUTVA] \textit{Stable Unit Treatment Values Assumption}.
    \item[RD] \textit{Risk Difference}, diferencia de riesgo.
    \item[RR] Riesgo relativo.
    \item[OR] \textit{Odds Ratio}.
    \item[ATE] \textit{Average Treatment Effect}, efecto medio del tratamiento.
    \item[IPTW] \textit{Inverse Probability of Treatment Weighting}.
    \item[HT] Horvitz-Thompson.
    \item[AIPW] \textit{Augmented Inverse Probability Weighting}.
\end{description}
\endinput


% -------------------------------------------------------------------
% BACKMATTER
% -------------------------------------------------------------------

\backmatter % Desactiva la numeración de los capítulos
\pdfbookmark[-1]{Referencias}{BM-Referencias}

% BIBLIOGRAFÍA
%-------------------------------------------------------------------

\bibliographystyle{abbrv}
% Otros estilos bibliográficos: alpha-es, amsalpha, plain
\begin{small}
  \bibliography{library.bib}
\end{small}

\end{document}
